%! Summary Statistics for a Quantitative Variable — Unit 1, Lesson 1.5 — Lesson Plan
% GRADUAL-RELEASE plan (warm-up / guided notes and practice / debrief / close and start homework).
% THERE IS NO GROUP ACTIVITY IN THIS COURSE — the guided notes carry the release: I do (two
% sections) -> we do (guided practice). THERE IS NO INDEPENDENT PRACTICE SET — the homework is
% the individual practice, and the period ends with students starting it. Teacher-facing. Every box is
% `skillbox` (breakable) with a \boxguard ahead of it; this course has no `fixedskillbox`,
% no tier boxes, and no exit ticket.
\documentclass[10pt]{article}
\usepackage{apstats-article}
\usepackage{apstats-boxes}

\newcommand{\UnitNumberName}{Unit 1: Exploring One-Variable Data \quad}
\newcommand{\LessonNumberName}{Lesson 1.5: Summary Statistics for a Quantitative Variable}

\begin{document}
\begin{center}
    {\Large\bfseries \CourseName \\
    {\small\bfseries \UnitNumberName \LessonNumberName}}
\end{center}

\begin{tcolorbox}[colback=sky, colframe=navy, boxrule=0.9pt, arc=2mm,
    left=2mm, right=2mm, top=1.5mm, bottom=1.5mm]
    \textbf{Primary Objective:} Students will calculate measures of center (mean, median),
    position (quartiles, percentiles), and variability (range, IQR, standard deviation) for a
    quantitative variable, and will \emph{justify} which pair of summaries to report from the
    shape of the distribution and the presence of an outlier.
    \\[2pt]
    \textbf{CED alignment:} Topic 1.7 \quad$\cdot$\quad Big Idea: UNC (Uncertainty)
    \quad$\cdot$\quad Skills 2.C, 4.B \quad$\cdot$\quad LO UNC-1.I, UNC-1.J, UNC-1.K
    \quad$\cdot$\quad EK UNC-1.I.1--1.I.5, 1.J.1--1.J.4, 1.K.1--1.K.2
    \\[2pt]
    \textbf{Lesson model:} \emph{Gradual release --- I do, we do, you do --- all of it inside
    the guided notes.} There is no group activity. The notes name the vocabulary as they teach
    it and deliver the instruction in two sections; the guided practice is worked together with
    students holding the pen; \textbf{the homework is the individual practice}, and students
    start it in class while you circulate. The whole-class debrief is \emph{spoken} ---
    there is no transfer set and no reflection box on the page.
    \\[2pt]
    \textbf{Note on arithmetic:} every data set in this lesson is built to divide cleanly, and
    every display is \textbf{pre-drawn}. Students never sketch a plot; the only ``construction''
    is filling a table. Keep calculators in the bag except for one thing --- the
    \emph{standard deviation}, which the notes compute by hand exactly once (five values) to
    show what it measures and then hand to technology in plain words. Every other number is
    faster with a pencil, and a calculator turns the paycheck arithmetic into a button hunt.
\end{tcolorbox}

\renewcommand{\arraystretch}{1.4}
\boxguard
\begin{skillbox}[Learning Targets \& Key Understandings]{goldbox}
\begin{tabularx}{\linewidth}{@{}X|X@{}}
\textbf{Learning targets (student language)} & \textbf{Key understandings (the ``why'')} \\[2pt]
\hline
\vspace{-6pt}
\begin{itemize}[leftmargin=1.1em, nosep, after=\vspace{2pt}]
  \item I can compute the mean and the median of a set of values, and find the quartiles by
        splitting the ordered list.
  \item I can show how one unusual value drags one summary a long way and leaves the other
        where it was.
  \item I can put a single number on how spread out a set of values is --- and say what that
        number can be fooled by.
  \item I can choose which pair of summaries to report and defend the choice from the shape.
\end{itemize}
&
\vspace{-6pt}
\begin{itemize}[leftmargin=1.1em, nosep, after=\vspace{2pt}]
  \item A \textbf{statistic} is a numerical summary of \emph{sample} data (EK UNC-1.I.1). The
        mean sums every value; the median only asks which value sits in the middle
        (EK UNC-1.I.2, 1.I.3). Quartiles and percentiles locate a value in the pack
        (EK UNC-1.I.4, 1.I.5).
  \item \textbf{Mean, standard deviation, and range are nonresistant}; \textbf{median and IQR
        are resistant} (EK UNC-1.K.2). That single fact is what makes the choice of summary a
        defensible decision rather than a preference.
  \item \textbf{Target misconception:} that the mean is simply ``the right answer'' for center
        and the median is a shortcut. The AP framing is the reverse --- both are correct, and
        the \emph{shape} and the presence of an outlier pick between them (Skill 4.B). The
        crux that surfaces it is the second half of notes section 2, \emph{One Value Changes}.
  \item \textbf{Second misconception:} that range is an \emph{interval} (``from 27 to 51'').
        Range is one number, $\text{max} - \text{min}$ (EK UNC-1.J.2).
\end{itemize}
\end{tabularx}
\end{skillbox}

\boxguard
\begin{skillbox}[{Vocabulary, Concepts, \& Theorems --- taught directly in the guided notes}]{greenbox}
\begin{minipage}[t]{0.48\linewidth}
    \begin{itemize}
        \item \textbf{Statistic} --- a numerical summary of sample data (EK UNC-1.I.1)
        \item \textbf{Mean} $\bar x = \sum x_i / n$ --- sum of the values over the number of
              values (EK UNC-1.I.2)
        \item \textbf{Median} --- the middle of the ordered data; average the two middle values
              when the count is even (EK UNC-1.I.3)
        \item \textbf{$Q_1$, $Q_3$} --- medians of the lower and upper halves; they bound the
              middle 50\% (EK UNC-1.I.4)
        \item \textbf{$p$th percentile} --- the value with $p\%$ of the data at or below it
              (EK UNC-1.I.5)
    \end{itemize}
\end{minipage}
\hfill
\begin{minipage}[t]{0.48\linewidth}
    \begin{itemize}
        \item \textbf{Range} --- $\text{maximum} - \text{minimum}$, a single number
              (EK UNC-1.J.2)
        \item \textbf{Interquartile range (IQR)} --- $Q_3 - Q_1$ (EK UNC-1.J.2)
        \item \textbf{Standard deviation} $s_x = \sqrt{\frac{1}{n-1}\sum (x_i - \bar x)^2}$,
              and its square the \textbf{variance} $s_x^2$ (EK UNC-1.J.3); changing units
              rescales every statistic (EK UNC-1.J.4)
        \item \textbf{Resistant / nonresistant} --- median and IQR vs.\ mean, standard
              deviation, and range (EK UNC-1.K.2); the two outlier rules (EK UNC-1.K.1)
        \item \textbf{Carried forward} --- shape, center, spread, unusual features, skew, and
              outlier (1.4); dotplots (1.3)
    \end{itemize}
\end{minipage}
\end{skillbox}

% A tabularx never splits, so guard generously ahead of it.
\boxguard[30]
\begin{skillbox}[Lesson at a Glance --- \MeetingLength]{sky}
\renewcommand{\arraystretch}{1.25}
\begin{tabularx}{\linewidth}{@{}l c X X@{}}
\textbf{Phase} & \textbf{Min} & \textbf{Students} & \textbf{Teacher} \\
\hline
Warm-Up & 5 & A shape from a dotplot, an average, and one ordered list cut into quarters & Debrief item~3 aloud; leave the split list on the board all period \\
Guided Notes \& Practice & 34 & Fill the vocabulary box and notes 1--2; then \emph{The Quiz Scores} together, pen in their hand & \textbf{I do} notes 1--2 (20) $\to$ \textbf{we do} the Guided Practice box (14) \\
Debrief & 8 & Pens down, papers up --- answer aloud, nothing to fill in & Open on the almost-right crux answer (``the median is more accurate''); put \$1{,}000 $\to$ \$750 and \$750 $\to$ \$750 side by side \\
Close \& Assign & 8 & \textbf{Start the homework in class}, alone and silent & Announce the paper homework, circulate the first minutes for your formative read, preview Lesson~1.6 \\
\end{tabularx}
\end{skillbox}

\boxguard[20]
\begin{skillbox}[Warm-Up --- Activate Prior Knowledge (5 min)]{sky}
\begin{minipage}[t]{0.48\linewidth}
    \textbf{The three items and what each seeds}
    \begin{itemize}
        \item \textbf{1.} Describe a right-skewed dotplot of television hours. \textbf{Spirals}
              to Lesson 1.4 --- and puts one far-out value in front of them minutes before
              section~1 needs it.
        \item \textbf{2.} Average six numbers; order them and find the middle. \textbf{Seeds:}
              the exact arithmetic of the paychecks, done once on easy numbers where the
              average and the middle happen to agree.
        \item \textbf{3.} Split an ordered list of seven, then split each half again.
              \textbf{Seeds:} the quartile move --- students do it before it has a name, and
              section~1 names it twenty minutes later.
    \end{itemize}
\end{minipage}
\hfill
\begin{minipage}[t]{0.48\linewidth}
    \textbf{Running it}
    \begin{itemize}
        \item \textbf{Debrief item~3 aloud} and \textbf{leave the board work up all period}:
              the seven numbers with three vertical strokes through them (lower half $|$
              middle $|$ upper half), 5, 9, and 15 circled. Section~1 attaches $Q_1$, the
              median, and $Q_3$ to that picture rather than to a definition supplied cold.
        \item Item~2 is a set where the average \emph{equals} the middle value. Nobody
              needs to notice yet --- the paychecks are what make that fact interesting.
        \item Item~1 is a 45-second spiral; do not let a shape debate expand. ``Skewed right,
              one student far out at 14'' is the whole answer.
    \end{itemize}
\end{minipage}
\end{skillbox}

\boxguard
\begin{skillbox}[Guided Notes \& Practice --- I do, then we do (34 min)]{sky}
\textbf{This is the whole lesson.} There is no group activity, and there is no independent
practice set on this page: \textbf{the homework is the individual practice}, started in class.
\begin{multicols}{2}
\textbf{I do (about 20 min).} Two sections, each in two moves; students fill the blanks as each
idea is named, and the vocabulary box is filled \emph{as you go}, never in advance.

\emph{Section 1} opens on \emph{statistic} and the nine paychecks: the mean in the \texttt{work}
block (\$9{,}000 over 9), then the ordered list and the fifth value. Land the line
\emph{nobody at this company takes home \$1{,}000} and say out loud that both numbers are
correct and section~2 decides which one to report --- do not resolve it yet. Five minutes. Its
second half, \emph{Quartiles and Percentiles}, is Route A: twelve times, the even-count median,
and the split move from the warm-up board with its names attached --- $Q_1 = 28.5$,
$Q_3 = 31.5$. Percentile gets one sentence and one interpretation in context. Five minutes.

\textbf{Section 2 is the lesson.} Its first half names the three spread numbers: range and IQR
in a \texttt{work} block --- say \emph{one number, not an interval} while you write $33 - 27$
--- then the standard deviation on the five-package route, the only hand computation of $s_x$
this course asks for. Show the table, the formula, and $\sqrt{10}$, then say plainly that the
calculator does this from now on, and that changing units rescales every statistic. Five
minutes. Its second half, \emph{One Value Changes}, is the crux: swap the \$3{,}000 for \$750
in the \texttt{work} block, write \$1{,}000 $\to$ \$750 beside \$750 $\to$ \$750, and make
someone say \emph{why} one of them did not care. Then Routes B and C: fill the table, and
let the room see range 24 = range 24 on the strength of one package while the IQR stays 3.
Land \emph{resistant} and \emph{nonresistant} on that table, then the reporting rule, then the
two outlier rules in one breath with $1.5\times\text{IQR}$ flagged for next lesson. Five
minutes.

\columnbreak
\textbf{We do (about 14 min).} The Guided Practice box, \emph{The Quiz Scores} --- eleven
scores so the median is a single value and each half holds exactly five. Students hold the pen;
you hold the questions: \emph{``Which position in the list did you change? Is that position
anywhere near the middle?''} \quad \emph{``The IQR only looks at two numbers. Did either of them
change?''} Part (c) is the crux transferred and the part to protect when time is short; part
(d) is Skill 4.B in one sentence.

\textbf{The pen must actually change hands.} Put the five-number table on them cold, take a
hand count on \emph{does the median move?} in part (c) before anyone speaks, and make the room
defend both answers before the \texttt{work} arithmetic settles it.

Circulate during Guided Practice and demand the reason, not the verdict:
\begin{itemize}[leftmargin=1.1em, nosep]
  \item \textbf{Part (a):} ``Read me the five below the median. Which one is their middle?''
  \item \textbf{Part (b), the interval trap:} ``I asked for one number. Which subtraction gives
        it?''
  \item \textbf{Part (c), the crux:} ``Four verdicts is not an answer. What rule sorts them
        into two piles?''
\end{itemize}
While circulating, pick the papers to put up in the debrief --- one whose part (c) says the
median ``is more accurate'' (almost right, wrong reason), and one that names positions.
\textbf{Your formative read comes twice}: here, and again in the last eight minutes as students
start the homework.
\end{multicols}
\end{skillbox}

\boxguard
\begin{skillbox}[Debrief --- whole class (8 min)]{sky}
Pens down, papers up. \textbf{This debrief is spoken} --- there is no box at the end of the
notes to fill in, so it lives entirely on the board and in the room.
\begin{multicols}{2}
\begin{enumerate}[leftmargin=1.3em, nosep, itemsep=3pt]
  \item \textbf{Open on the almost-right answer.} Put up the Guided Practice paper that says
        \emph{``the median is more accurate.''} It is not more accurate --- the mean of the
        corrected scores really is 15. It is \emph{resistant}: it only asks which score is
        sixth. Write \$1{,}000 $\to$ \$750 and \$750 $\to$ \$750 beside it and make a student
        say which of the two arithmetic procedures ever touched the \$3{,}000.
  \item Put up the paper that named positions and say what it did right: it explained the
        \emph{mechanism}, and the mechanism is the rule (EK UNC-1.K.2).
  \item Take the interval trap off part (b): circle a ``7'' and a ``12 to 19'' and ask which one
        you could add to another route's spread.
  \item Close on Skill 4.B: \emph{``Mean and median are both correct. The shape --- and whether
        one strange value is sitting in it --- is what decides which one you are allowed to
        quote.''} Say it, write it, point at the paychecks.
\end{enumerate}
\columnbreak
\textbf{Two things to demand aloud.} First, \textbf{make a student name the person who earns the
mean} --- there is none --- and then the person who earns the median. Second, \textbf{make them
give a data set where mean and standard deviation are the right pair}, so nobody leaves
believing the median is always better; Route A is sitting on the page for exactly that purpose.

\textbf{Connect back to Lesson~1.4 if time allows:} a right-skewed dotplot and a mean well above
the median are the same fact read twice --- the tail is where the mean went.

\textbf{If the debrief runs short}, cut items 2 and 3 --- item~1 is the only one that cannot be
cut. \textbf{Do not let the debrief eat the last eight minutes} --- the homework start is not
optional padding, it is your second formative read.
\end{multicols}
\end{skillbox}

\boxguard
\begin{skillbox}[AP Practice --- assigned, not scored]{goldbox}
Fifteen households' daily water use, right-skewed with one home at 620 gallons, given as a
stem plot \emph{and} as computer output --- because that pairing is what the exam actually hands
students. Four multiple-choice items: \textbf{the crux as a multiple-choice item} (which pair to
report, and the reason is the shape); a resistance judgment (double the largest value and ask
which summary moves most); \textbf{the units item (EK UNC-1.J.4)}, where converting gallons to
litres multiplies mean, median, range, IQR, \emph{and} standard deviation alike; and a
percentile interpretation. Then one free response: name the shape, choose a center--spread pair
and \emph{justify the choice from the shape}, interpret $Q_1$ in context, and rebut a neighbour
who quotes the mean as ``typical.''

\textbf{Part (a) is the spiral} --- it reaches straight back to Lesson 1.4, and part (b) is
Skill~4.B in the exact form the exam asks it: the choice earns nothing on its own; the choice
justified by the skew and the one large value earns everything.

\textbf{Not scored --- the cover's score column reads \textbf{NA} for this page.} The water
record appears nowhere else in the packet, so what comes back in Lesson~1.6 is your evidence
that the idea transferred to a data set students met on their own. Go over it at the start of
Lesson~1.6 and hold the AP standard out loud: \emph{in context or it does not count}. ``The
median is better'' earns nothing; ``the median, 200 gallons, because the distribution is skewed
right and the home using 620 gallons pulls the mean up to 240'' earns everything.
\end{skillbox}

\boxguard
\begin{skillbox}[Homework --- scored, due the first class after two study halls]{goldbox}
Five exercises spanning Topic 1.7 in fresh contexts: \textbf{1} mean and median for seven
commutes, and which one represents the office, \emph{given its shape}; \textbf{2} the resistance
crux restated --- one value changed, report what moved; \textbf{3} $Q_1$, $Q_3$, and IQR from an
ordered list of twelve, which forces the even-count rule three times; \textbf{4} a percentile
interpreted in one sentence; \textbf{5} the AP-style multiple-choice item --- mean \$95{,}000
against median \$58{,}000 --- with a one-sentence justification.

\textbf{Start it in the last eight minutes of the period, in class and alone.} \textbf{Scoring:}
10 points --- 2 per item. \textbf{Item~2 is where the misconception shows}: a student who
reports that the median changed has re-sorted the list carelessly, and a student who reports
that nothing but the mean changed has forgotten the range. \textbf{Item~5 carries the check:} the
mean salary against the median is the nine paychecks in a suit.

\textbf{Paper homework is the default for this course} --- assign this page unless you have
decided otherwise. \textbf{DeltaMath override (occasional):} on the days you assign DeltaMath
instead, it replaces this page, you strike through the score cell on the cover, and this page
stays in the packet as extra practice. Never assign both.
\end{skillbox}

\boxguard
\begin{skillbox}[Watch For (while circulating)]{redbox}
\begin{itemize}[leftmargin=1.2em, nosep]
  \item \textbf{``The mean is the real answer; the median is a shortcut''} (notes~1 and~2,
        GP~(c)--(d), HW~1, AP~MC1 and~(b) --- the target misconception). Probe: \emph{``Name the
        person who earns the mean. Now name the person who earns the median.''} The
        almost-right version --- \emph{``the median is more accurate''} --- opens the debrief.
  \item \textbf{Median found without ordering the list} (notes~1, HW~1). The paycheck list is
        printed out of order on purpose. Probe: \emph{``Read me your list from smallest to
        largest first.''}
  \item \textbf{Range reported as an interval} --- ``the range is 27 to 51'' (notes~2, GP~(b),
        HW~2). Probe: \emph{``I asked for one number. Which subtraction gives it?''}
  \item \textbf{Even-count median taken as one of the two middle values} (notes~1, HW~3): with
        twelve times they must average. Probe: \emph{``Six below and six above --- so which
        single one is in the middle? Neither. What now?''}
  \item \textbf{IQR computed as $Q_1$ to $Q_3$, or the halves split wrongly} (notes~1, GP~(a),
        HW~3). Probe: \emph{``Go back to the warm-up on the board. How many numbers were in each
        half?''}
  \item \textbf{Standard deviation assumed resistant} (notes~2, AP~MC2, HW~5): it is built from
        the mean, so it is not. Probe: \emph{``What is inside the formula? Then how could it
        survive a \$3{,}000 paycheck?''}
  \item \textbf{Units ignored} (AP~MC3): a summary statistic carries the units of the variable.
        Probe: \emph{``Seven what?''}
\end{itemize}
Cold-call prompts: \emph{``Give me a data set where the mean and the median are far apart.''}
\quad \emph{``Which summary would you hide if you were the boss?''}
\end{skillbox}

\boxguard
\begin{skillbox}[Close \& Assign (8 min)]{goldbox}
\textbf{Students start the homework here, in class, alone and silent --- this is the individual
practice, and the first minutes of it are your best formative read.} Hand the launch first ---
five exercises on paper, scored, due the first class after two study halls --- and point
out that the AP Practice page before it is theirs to work and bring to review. Circulate:
\textbf{item~2 is the column that tells you whether section~2 landed} --- five values, one
changed to 32, and the median does not move. Sort what you see into three piles as you walk
--- said the mean and range moved and named \emph{why} the median did not (ready); got the three
numbers right with no reason (ready, needs the language); reported that the median moved, or
that only the mean did (the misconception survived) --- and open Lesson~1.6 by reteaching the
second half of section~2 if that third pile ran past a third of the room. Then close on what
changed: \emph{``You walked in thinking there was one right number for `typical.' You walk out
knowing there are two, that they can be \$250 apart, and that the shape of the picture is what
tells you which one you are allowed to quote.''} \textbf{Preview:} Lesson~1.6, \emph{Boxplots
and Comparing Distributions} --- the five numbers they just computed become a picture, and that
picture is what lets us compare two groups side by side. The $1.5\times\text{IQR}$ rule is named
today and checked once on Route C; it does its real job next lesson, where the whiskers stop.
\end{skillbox}

% ── Teacher notes — one per component, in packet order (four: there is no activity) ──
\begin{teachernote}[Warm-Up]
Four minutes, five at the outside. Item~1 is a 45-second spiral; do not let a shape debate
expand. Item~2 is mechanical and its punchline is invisible on purpose --- the average and the
middle value both come out to 8, which is exactly the situation the paychecks are about to
break. \textbf{Item~3 is the one to debrief aloud, and its board work stays up all period}:
seven numbers, the middle circled, then each half's middle circled. You will point at that board
twice --- once when section~1 names $Q_1$ and $Q_3$ on Route A, and once during the Guided
Practice table. Use the students' own words here (``the middle one,'' ``the middle of each
half''); the notes attach \emph{median} and \emph{quartile} to them twenty minutes later.
\end{teachernote}

\begin{teachernote}[Guided Notes \& Practice]
\textbf{This one document is the lesson --- pace it 20 / 14, then 8 for the debrief, then 8 for
the homework start.} Section~1 is the machinery and must move: five minutes on the paychecks,
five on Route A. Compute the mean in the \texttt{work} block live and make a student order the
nine before you take the fifth --- the list is printed out of order on purpose. The line to land
in section~1 is that nobody earns \$1{,}000; the line to \emph{withhold} is which number to
report. Say plainly that section~2 decides it.

\textbf{The first half of section 2 is where the range becomes one number.} Write $33 - 27$ and
say ``six --- not twenty-seven to thirty-three'' before anyone writes an interval. The standard
deviation gets one hand computation, the formula, and the sentence \emph{the calculator does
this from now on}; do not let it expand past five minutes, because nothing in this packet asks
for a second one by hand.

\textbf{Spend the time in the second half of section 2}, which carries the misconception. Swap
the \$3{,}000 for \$750 in the \texttt{work} block and write the two arrows side by side before
saying \emph{resistant}. Then fill the routes table and let the room see 24 against 24 --- the
range cannot tell Route C from Route B --- before you point at the IQR column. The word
\emph{resistant} lands on that table, and the reporting rule lands on the word. Do not resolve
the tension too fast: the sentence to land is that both measures of center are correct and the
shape picks.

\textbf{Guided Practice is where the pen changes hands.} \emph{The Quiz Scores} is eleven values
so the table is clean; part (c) is the crux transferred and part (d) is the reporting rule in
one sentence. If you only get through (a) and (c), that is the right trade.

\textbf{There is no independent practice set on this page, and that is deliberate --- the
homework is the individual practice.} Guided Practice is the last thing on the page; the period
ends with students starting the homework in front of you, which is where your second formative
read comes from. Item~2 is the one to watch: five values, one changed to 32, and the median stays
10. Sort as you circulate --- named why the median did not move (ready); three numbers right with
no reason (ready, needs the language); reported the median moved (the misconception survived).
If the third pile is more than a third of the room, \textbf{open Lesson~1.6 by reteaching the
second half of section~2}, and say so plainly.

\textbf{Debrief (8 min), spoken.} Open on the almost-right answer --- \emph{``the median is more
accurate''} --- and replace \emph{accurate} with \emph{resistant} on the board; it is the one
move that cannot be cut. \textbf{Do not let the debrief eat the last eight minutes} --- the
homework start is not optional padding.
\end{teachernote}

\begin{teachernote}[AP Practice]
\textbf{This page is not required and not scored} --- the graded homework is the section behind
it at the back of the packet. Work it as AP-format reps and go over it at the start of
Lesson~1.6. \textbf{MC item~1 is the crux in exam clothing}: distractor (A) --- ``mean and
standard deviation, because they use every value'' --- is the target misconception stated as a
virtue, and (D) pairs a resistant center with a nonresistant spread. If more than a handful pick
(A), reopen the paychecks. \textbf{MC item~3 is the other one worth reading closely}: converting
gallons to litres multiplies \emph{every} one of these statistics by the conversion factor, and
distractor (B) --- ``the standard deviation is unchanged'' --- is the intuition that a spread
measure is somehow unit-free. It is not (EK UNC-1.J.4), and the same confusion returns in Unit~4
with linear transformations of random variables. On the free response, part (b) is the
highest-value part: it forces students to justify a \emph{choice} of summary from the shape,
which is Skill~4.B in the form the exam actually asks it, and an unjustified ``median and IQR''
earns nothing.
\end{teachernote}

\begin{teachernote}[Homework]
\textbf{Scored, 10 points (2 per item), due the first class after two study halls.} The eight
minutes at the end of the period are a \emph{start}, not a completion: put students on 1, 2,
and~5 while you can still catch a wrong start, and let 3 and~4 go home. Item~2 is the first one
to read when scoring, because it is the crux stripped of context: mean $10\to14$, range
$4\to24$, median unchanged at 10. Item~3 is the cheapest diagnostic in the set --- a student who
reports $Q_1=30$ or $Q_3=85$ has taken one of the two middle values instead of averaging them,
and that error will wreck every boxplot in Lesson~1.6. \textbf{When you score the page, sort
item~5 into three piles} --- chose (B) and named the skew (ready); chose (B) and said only
``because of outliers'' (ready, needs the language); chose (A) or (C) (the misconception
survived). If more than a third land in the last pile, reopen the nine paychecks in the
Lesson~1.6 warm-up debrief rather than letting it ride. Hold the context line while scoring:
item~4 must name newborn girls and grams, not merely ``90\% are below.''

\textbf{Paper is the default.} This page is what gets assigned unless you decide otherwise on a
given day. On the occasions you assign DeltaMath in its place, strike the score cell on the cover
and tell students this page is now optional practice. Do not assign both.
\end{teachernote}
\end{document}
